% !TeX program = xelatex
\documentclass[12pt,a4paper]{extarticle}

\usepackage{fontspec}
\usepackage{amssymb}
\usepackage{polyglossia}
\setdefaultlanguage{russian}
\setmainfont{DejaVu Serif}

\usepackage{geometry}
\geometry{
    left=3cm,
    right=1.5cm,
    top=2cm,
    bottom=2cm
}

\usepackage{setspace}
\onehalfspacing

\usepackage{indentfirst}
\usepackage{tikz}
\usepackage{makecell}
\usepackage{amsmath}
\usepackage{graphicx}
\usepackage{hyperref}
\usepackage{float}
\usepackage{multirow}
\usepackage{array}

\usepackage{caption}
\captionsetup{
    font={small,it}
}

\hypersetup{
    colorlinks=false,
    pdftitle={Домашняя работа 1}
}

\renewcommand{\arraystretch}{1.2}

\begin{document}

% Титульный лист
\begin{titlepage}
    \begin{center}
        Федеральное государственное автономное образовательное учреждение высшего образования «Национальный исследовательский университет ИТМО»\\
        Факультет программной инженерии и компьютерной техники

        \vfill

        \textbf{Домашняя работа №1}

        по предмету «Дискретная математика»

        на тему:

        \textbf{«Название темы»}

        Вариант №14.
        \vfill
    \end{center}

    \begin{flushright}
        Выполнил:\\
        Вишняков Александр Сергеевич\\
        группа Р3110
        
        \vspace{0.5cm}

        Проверил:\\
        Поляков В.И.
    \end{flushright}

    \vfill

    \begin{center}
        Санкт-Петербург, 2026
    \end{center}
\end{titlepage}

\tableofcontents
\clearpage

\section*{Задания}
\addcontentsline{toc}{section}{Задания}

\subsection*{1. Представление чисел с фиксированной и плавающей запятой в различных форматах}
\addcontentsline{toc}{subsection}{Представление чисел с фиксированной и плавающей запятой в различных форматах}

\begin{enumerate}
    \item Заданное число $A$ представить в виде двоично-кодированного
    десятичного числа:
    \begin{enumerate}
        \item в упакованном формате (BCD);
        \item в неупакованном формате (ASCII).
    \end{enumerate}

    \item Заданные числа $A$ и $-A$ представить в форме с фиксированной
    запятой.

    \item Заданные числа $A$ и $B$ представить в форме с плавающей
    запятой в формате Ф1.

    \item Заданные числа $A$ и $B$ представить в форме с плавающей
    запятой в формате Ф2.

    \item Заданные числа $A$ и $B$ представить в форме с плавающей
    запятой в формате Ф3.

    \item Найти значения чисел $Y$ и $Z$ по их заданным
    шестнадцатеричным представлениям $R$ и $S$ в форме с плавающей
    запятой в формате Ф1.

    \item Найти значения чисел $V$ и $W$ по их заданным
    шестнадцатеричным представлениям $R$ и $S$ в форме с плавающей
    запятой в формате Ф2.

    \item Найти значения чисел $T$ и $Q$ по их заданным
    шестнадцатеричным представлениям $R$ и $S$ в форме с плавающей
    запятой в формате Ф3.
\end{enumerate}

\noindent\textbf{Замечание.}
При выполнении пп.~3--5 задания для дробного числа $B$ в целях
увеличения точности его представления произвести симметричное
округление мантиссы.

\subsection*{Заданные числа}
\addcontentsline{toc}{subsection}{Заданные числа}

\begin{table}[htbp]
    \centering
    \caption{Заданные значения из варианта 14}
    \begin{tabular}{|c|c|c|c|}
		$A$ & $B$ & $R$ & $S$ \\
		\hline
		$1400$ & $0.05$ & $C1E08000$ & $3EA00000$ \\
		
    \end{tabular}
\end{table}


\subsection*{2. Сложение целых чисел}
\addcontentsline{toc}{subsection}{Сложение целых чисел}

\begin{enumerate}
    \item Для заданных чисел $A$ и $B$ выполнить операцию знакового
    сложения со всеми комбинациями знаков операндов. Для каждого
    примера:
    \begin{enumerate}
        \item проставить межразрядные переносы, возникающие при сложении;
        \item дать знаковую интерпретацию (Зн) операндов и результатов.
        При получении отрицательного результата предварительно
        преобразовать его из дополнительного кода в прямой;
        \item дать беззнаковую интерпретацию (БзИ) операндов и
        результатов. При получении неверного результата пояснить
        причину его возникновения;
        \item показать значения арифметических флагов.
    \end{enumerate}

    \item Сохранив значение первого операнда $A$, выбрать такое
    значение $B$, чтобы в операции сложения с одинаковыми знаками
    имел место особый случай переполнения формата. Выполнить два
    примера, иллюстрирующие эти случаи. Для каждого из них выполнить
    пункты а--d.

    \item Сохранив операнд $B$, подобрать такое значение операнда $A$,
    чтобы при сложении положительных операндов имело место
    переполнение формата, а при сложении отрицательных операндов
    результат был бы корректен. Выполнить два примера,
    иллюстрирующие этот случай. Для каждого из них выполнить
    пункты а--d.
\end{enumerate}

\subsection*{Заданные числа}
\addcontentsline{toc}{subsection}{Заданные числа}

\begin{table}[H]
    \centering
    \caption{Заданные значения из варианта 14}
    \begin{tabular}{|c|c|}
		$A$ & $B$ \\
		\hline
		$109$ & $17$  \\
		
    \end{tabular}
\end{table}



\subsection*{3. Вычитание целых чисел}
\addcontentsline{toc}{subsection}{Вычитание целых чисел}

\begin{enumerate}
    \item Для заданных чисел $A$ и $B$ выполнить операцию знакового
    вычитания со всеми комбинациями знаков операндов. Для каждого
    примера:
    \begin{enumerate}
        \item проставить межразрядные заёмы, возникающие при вычитании;
        \item дать знаковую интерпретацию (Зн) операндов и результатов.
        При получении отрицательного результата предварительно
        преобразовать его из дополнительного кода в прямой;
        \item дать беззнаковую интерпретацию (БзИ) операндов и
        результатов. При получении неверного результата пояснить
        причину его возникновения;
        \item показать значения арифметических флагов.
    \end{enumerate}

    \item Сохранив значение первого операнда $A$, выбрать такое
    значение $B$, чтобы в операции вычитания с разными знаками имел
    место особый случай переполнения формата. Выполнить два примера,
    иллюстрирующих эти случаи. Для каждого из них выполнить пункты а--d.

    \item Сохранив операнд $B$, подобрать такое значение операнда $A$,
    чтобы при вычитании отрицательного $B$ из положительного $A$ имело
    место переполнение формата, а при вычитании положительного $B$ из
    отрицательного $A$ результат был бы корректен. Выполнить два
    примера, иллюстрирующих этот случай. Для каждого из них выполнить
    пункты а--d.
\end{enumerate}

\subsection*{Заданные числа}
\addcontentsline{toc}{subsection}{Заданные числа}

\begin{table}[H]
    \centering
    \caption{Заданные значения из варианта 14}
    \begin{tabular}{|c|c|}
		$A$ & $B$ \\
		\hline
		$109$ & $17$  \\
		
    \end{tabular}
\end{table}


\section{Представление чисел с фиксированной и плавающей запятой в различных форматах}

\subsection{Представление числа в упакованном и неупакованном форматах}
Получим двоичное представление каждого разряда числа $A = 1400$, дополнив каждое число до 4 бит,
получая тем самым двоичные тетрады:

$1_{10} = 0001_{2}$

$4_{10} = 0100_{2}$

$0_{10} = 0000_{2}$

$0_{10} = 0000_{2}$

В упакованном формате $BCD$ каждый байт кодирует два десятичных числа (две тетрады).
Представление числа в упакованном виде представлено в таблице~\ref{tab:a-in-bcd}.

\begin{table}[H]
    \centering
    \caption{Представление числа в упакованном формате}
    \label{tab:a-in-bcd}
    \begin{tabular}{|c|c|}
        \hline
		$0001.0100$ & $0000.0000$ \\
		\hline
    \end{tabular}
\end{table}

В неупакованном формате $ASCII$ каждый байт кодирует одну десятичную цифру (одну тетраду),
а первая тетрада каждого байта имеет стандартное значение $0011$.
Представление числа в неупаковааном виде представлено в таблице~\ref{tab:a-in-ascii}.

\begin{table}[H]
    \centering
    \caption{Представление числа в неупакованном формате}
    \label{tab:a-in-ascii}
    \begin{tabular}{|c|c|c|c|}
        \hline
		$0011.0001$ & $0011.0100$ & $0011.0000$ & $0011.0000$\\
		\hline
    \end{tabular}
\end{table}

\subsection{Представление числа в форме с фиксированной запятой}

\subsubsection{Представление числа A}
$A = 1400_{10} = 10101111000_{2}$

Число $A$ занимает 11 бит, а также нужен 1 бит под знак числа. Поэтому возьмём $n = 12$.
Так как $A > 0$, то его представление будет в форме с фиксированной запятой будет неизменным,
только добавится первый бит справа -- знаковый (равный $0$).

$A = 0.10101111000_{2}$

\subsubsection{Представление числа -A}
Аналогично предыдущему примеру, возьмём $n = 12$. Так как число $-A = -1400 < 0$, то
его двоичному представлению будет соответствовать форма $2^n - |A|$.

$-A = 2^n - 1400 = 4096 - 1400 = 2696_{10} = 1.01010001000_{2}$


\subsection{Представление чисел в формате Ф1}

\subsubsection{Представление числа A}
$A = 1400_{10} = 0101 0111 1000_{2} = 578_{16}$

В формате Ф1 используется основание системы счисления $S = 16$. Необходимо представить число в виде
$A = (-1)^{sign A} \cdot M_{A} \cdot S^{P_{A}}$.

$A = (0.578)_{16} \cdot 16^{3}$

Таким образом, мантисса $M_{A} = 0.578$, а порядок $P_{A} = 3$. В формате Ф1 используется смещение
порядка $d = 64$. Так, получим характеристику числа $X_{A}=P_{A}+d=3+64=67_{10}=1000011_{2}$. Записав знаковый 
бит, характеристику и мантиссу, дополненную до 6-ти тетрад, получим 32-битное представление числа
в формате Ф1. 

\begin{figure}[H]
    \centering
    \begin{tikzpicture}[
        bit/.style={font=\large, minimum width=0.48cm,
                    minimum height=0.65cm, inner sep=0pt},
        every node/.style={inner sep=0pt}
    ]

    \def\bits{
        0,1,0,0,0,0,1,1,
        0,1,0,1,0,1,1,1,
        1,0,0,0,0,0,0,0,
        0,0,0,0,0,0,0,0
    }

    \draw (0,0) rectangle (15.36,0.7);
    \draw (0,0) rectangle (0.48,0.7);
    \draw (0,0) rectangle (3.84,0.7);

    \foreach \b [count=\i from 0] in \bits {
        \node[bit] at ({0.24+0.48*\i},0.35) {\b};
    }

    \node[font=\small] at (0.24,1.05) {0};
    \node[font=\small] at (0.72,1.05) {1};
    \node[font=\small] at (3.58,1.05) {7};
    \node[font=\small] at (4.06,1.05) {8};
    \node[font=\small] at (15.12,1.05) {31};

    \draw (3.84,-0.05) -- (3.84,-0.25);
    \draw (3.84,-0.25) -- (5.76,-0.25);
    \draw (5.76,-0.05) -- (5.76,-0.25);

    \node[font=\small] at (4.9,-0.65) {5};

    \draw (5.76,-0.05) -- (5.76,-0.25);
    \draw (5.76,-0.25) -- (7.68,-0.25);
    \draw (7.68,-0.05) -- (7.68,-0.25);

    \node[font=\small] at (6.82,-0.65) {7};

    \draw (7.68,-0.05) -- (7.68,-0.25);
    \draw (7.68,-0.25) -- (9.6,-0.25);
    \draw (9.6,-0.05) -- (9.6,-0.25);

    \node[font=\small] at (8.74,-0.65) {8};

    \draw (0.48,-1.05) -- (0.48,-1.25);
    \draw (0.48,-1.25) -- (3.84,-1.25);
    \draw (3.84,-1.05) -- (3.84,-1.25);

    \node[font=\small] at (2.16,-1.65) {характеристика};

    \draw (3.84,-1.05) -- (3.84,-1.25);
    \draw (3.84,-1.25) -- (15.12,-1.25);
    \draw (15.12,-1.05) -- (15.12,-1.25);

    \node[font=\small] at (9.48,-1.65) {мантисса};

    \end{tikzpicture}
\end{figure}

\subsubsection{Представление числа B}

Переведём число $B = 0.05$ в шестнадцатеричную систему. Необходимо получить шесть значащих
разрядов, а также седьмой для округления числа.

\begin{table}[H]
    \centering
    \caption{Перевод числа В в шестнадцатеричную систему}
    \begin{tabular}{|c|c|c|}
        \hline
        \small Умножение & \small Целая часть & \small Дробная часть \\
        \hline
        $0.05 \cdot 16$ & \textbf{0} & $0.8$ \\
        $0.8 \cdot 16$ & \textbf{12=C} & $0.8$ \\
        $0.8 \cdot 16$ & \textbf{C} & $0.8$ \\
        $0.8 \cdot 16$ & \textbf{C} & $0.8$ \\
        $0.8 \cdot 16$ & \textbf{C} & $0.8$ \\
        $0.8 \cdot 16$ & \textbf{C} & $0.8$ \\
        $0.8 \cdot 16$ & \textbf{C} & $0.8$ \\
        $0.8 \cdot 16$ & \textbf{C} & $0.8$ \\
        \hline		
    \end{tabular}
\end{table}

Так как последний разряд $C=12_{10}=1100_{2}$ первым символом имеет 1, то для округления добавим единицу
к предыдущему разряду. Таким образом, $B = 0.0CCCCCD_{16}$.

Представим число в виде $B=(0.CCCCCD)_{16} \cdot 16^{-1}$. Тогда характеристика числа будет
$X_{A}=P_{A}+d=-1+64=63_{10}=0111111_{2}$. Запишем знаковый бит 0 и полученные значения в таблицу,
получив представление числа в формате Ф1.

\begin{figure}[H]
    \centering
    \begin{tikzpicture}[
        bit/.style={font=\large, minimum width=0.48cm,
                    minimum height=0.65cm, inner sep=0pt},
        every node/.style={inner sep=0pt}
    ]

    \def\bits{
        0,0,1,1,1,1,1,1,
        1,1,0,0,1,1,0,0,
        1,1,0,0,1,1,0,0,
        1,1,0,0,1,1,0,1
    }

    \draw (0,0) rectangle (15.36,0.7);
    \draw (0,0) rectangle (0.48,0.7);
    \draw (0,0) rectangle (3.84,0.7);

    \foreach \b [count=\i from 0] in \bits {
        \node[bit] at ({0.24+0.48*\i},0.35) {\b};
    }

    \node[font=\small] at (0.24,1.05) {0};
    \node[font=\small] at (0.72,1.05) {1};
    \node[font=\small] at (3.58,1.05) {7};
    \node[font=\small] at (4.06,1.05) {8};
    \node[font=\small] at (15.12,1.05) {31};

    \draw (3.84,-0.05) -- (3.84,-0.25);
    \draw (3.84,-0.25) -- (5.76,-0.25);
    \draw (5.76,-0.05) -- (5.76,-0.25);

    \node[font=\small] at (4.9,-0.65) {C};

    \draw (5.76,-0.05) -- (5.76,-0.25);
    \draw (5.76,-0.25) -- (7.68,-0.25);
    \draw (7.68,-0.05) -- (7.68,-0.25);

    \node[font=\small] at (6.82,-0.65) {C};

    \draw (7.68,-0.05) -- (7.68,-0.25);
    \draw (7.68,-0.25) -- (9.6,-0.25);
    \draw (9.6,-0.05) -- (9.6,-0.25);

    \node[font=\small] at (8.74,-0.65) {C};

    \draw (9.6,-0.05) -- (9.6,-0.25);
    \draw (9.6,-0.25) -- (11.52,-0.25);
    \draw (11.52,-0.05) -- (11.52,-0.25);

    \node[font=\small] at (10.66,-0.65) {C};

    \draw (11.52,-0.05) -- (11.52,-0.25);
    \draw (11.52,-0.25) -- (13.44,-0.25);
    \draw (13.44,-0.05) -- (13.44,-0.25);

    \node[font=\small] at (12.58,-0.65) {C};

    \draw (13.44,-0.05) -- (13.44,-0.25);
    \draw (13.44,-0.25) -- (15.36,-0.25);
    \draw (15.36,-0.05) -- (15.36,-0.25);

    \node[font=\small] at (14.5,-0.65) {D};

    \draw (0.48,-1.05) -- (0.48,-1.25);
    \draw (0.48,-1.25) -- (3.84,-1.25);
    \draw (3.84,-1.05) -- (3.84,-1.25);

    \node[font=\small] at (2.16,-1.65) {характеристика};

    \draw (3.84,-1.05) -- (3.84,-1.25);
    \draw (3.84,-1.25) -- (15.12,-1.25);
    \draw (15.12,-1.05) -- (15.12,-1.25);

    \node[font=\small] at (9.48,-1.65) {мантисса};

    \end{tikzpicture}
\end{figure}


\subsection{Представление чисел в формате Ф2}
\subsubsection{Представление числа А}
$A = 1400_{10} = 101 0111 1000_{2}$

В формате Ф2 используется основание системы $S=2$ и величина смещения порядка $d=128$.

$A = (0.10101111)_{2} \cdot 2^{11}$

$X_{A} = P_{A}+d=11+128=139_{10}=10001011_{2}$

Отметим также, что в формате Ф2 мантисса может быть представлена только в форме
нормализованного числа. Это значит, что первая цифра после запятой всегда будет
$1$, которая называется скрытым разрядом и не учитывается в итоговой записи.

\begin{figure}[H]
    \centering
    \begin{tikzpicture}[
        bit/.style={font=\large, minimum width=0.48cm,
                    minimum height=0.65cm, inner sep=0pt},
        every node/.style={inner sep=0pt}
    ]

    \def\bits{
        0,1,0,0,0,1,0,1,
        1,0,1,0,1,1,1,1,
        0,0,0,0,0,0,0,0,
        0,0,0,0,0,0,0,0
    }

    \draw (0,0) rectangle (15.36,0.7);
    \draw (0,0) rectangle (0.48,0.7);
    \draw (0,0) rectangle (4.32,0.7);

    \foreach \b [count=\i from 0] in \bits {
        \node[bit] at ({0.24+0.48*\i},0.35) {\b};
    }

    \node[font=\small] at (0.24,1.05) {31};
    \node[font=\small] at (0.77,1.05) {30};
    \node[font=\small] at (3.96,1.05) {23};
    \node[font=\small] at (4.66,1.05) {22};
    \node[font=\small] at (15.12,1.05) {0};


    \draw (0.48,-0.05) -- (0.48,-0.25);
    \draw (0.48,-0.25) -- (4.32,-0.25);
    \draw (4.32,-0.05) -- (4.32,-0.25);

    \node[font=\small] at (2.16,-0.65) {характеристика};

    \draw (4.32,-0.05) -- (4.32,-0.25);
    \draw (4.32,-0.25) -- (15.12,-0.25);
    \draw (15.12,-0.05) -- (15.12,-0.25);

    \node[font=\small] at (9.48,-0.65) {мантисса};

    \end{tikzpicture}
\end{figure}

\subsubsection{Представление числа В}
$B = 0.05_{10}=0.0CCCCCCC_{16}=0.0000 1100 1100 1100 1100 1100 1100 1100_{2}$

Округлив до 6-ти значимых тетрад, получим значение:

$B=0.0000 1100 1100 1100 1100 1100 1101_{2}$.

$B = (0.110011001100110011001101)_{2} \cdot 2^{-4}$

$X_{A} = -4 + 128 = 124_{10} = 01111100_{2}$

Так как в формате Ф2 характеристика занимает 8 бит, добавим 0 в начало числа.
Запишем знаковый бит 0 и полученные значения в таблицу, учитывая скрытый бит в мантиссе,
и получим представление в формате Ф2.

\begin{figure}[H]
    \centering
    \begin{tikzpicture}[
        bit/.style={font=\large, minimum width=0.48cm,
                    minimum height=0.65cm, inner sep=0pt},
        every node/.style={inner sep=0pt}
    ]

    \def\bits{
        0,0,1,1,1,1,1,0,
        0,1,0,0,1,1,0,0,
        1,1,0,0,1,1,0,0,
        1,1,0,0,1,1,0,1
    }

    \draw (0,0) rectangle (15.36,0.7);
    \draw (0,0) rectangle (0.48,0.7);
    \draw (0,0) rectangle (4.32,0.7);

    \foreach \b [count=\i from 0] in \bits {
        \node[bit] at ({0.24+0.48*\i},0.35) {\b};
    }

    \node[font=\small] at (0.24,1.05) {31};
    \node[font=\small] at (0.77,1.05) {30};
    \node[font=\small] at (3.96,1.05) {23};
    \node[font=\small] at (4.66,1.05) {22};
    \node[font=\small] at (15.12,1.05) {0};


    \draw (0.48,-0.05) -- (0.48,-0.25);
    \draw (0.48,-0.25) -- (4.32,-0.25);
    \draw (4.32,-0.05) -- (4.32,-0.25);

    \node[font=\small] at (2.16,-0.65) {характеристика};

    \draw (4.32,-0.05) -- (4.32,-0.25);
    \draw (4.32,-0.25) -- (15.12,-0.25);
    \draw (15.12,-0.05) -- (15.12,-0.25);

    \node[font=\small] at (9.48,-0.65) {мантисса};

    \end{tikzpicture}
\end{figure}


\subsection{Представление чисел в формате Ф3}
\subsubsection{Представление числа А}
$A = 1400_{10} = 101 0111 1000_{2}$

В формате Ф3 используется основание системы $S=2$ и величина смещения порядка $d=127$.
Также в этом формате запятая фиксируется после страшей единицы.

$A = (1.0101111)_{2} \cdot 2^{10}$

$X_{A} = P_{A}+d=10+127=137_{10}=10001001_{2}$

В этом формате скрытый разряд -- единица, после которой фиксируется запятая мантиссы.

\begin{figure}[H]
    \centering
    \begin{tikzpicture}[
        bit/.style={font=\large, minimum width=0.48cm,
                    minimum height=0.65cm, inner sep=0pt},
        every node/.style={inner sep=0pt}
    ]

    \def\bits{
        0,1,0,0,0,1,0,0,
        1,0,1,0,1,1,1,1,
        0,0,0,0,0,0,0,0,
        0,0,0,0,0,0,0,0
    }

    \draw (0,0) rectangle (15.36,0.7);
    \draw (0,0) rectangle (0.48,0.7);
    \draw (0,0) rectangle (4.32,0.7);

    \foreach \b [count=\i from 0] in \bits {
        \node[bit] at ({0.24+0.48*\i},0.35) {\b};
    }

    \node[font=\small] at (0.24,1.05) {31};
    \node[font=\small] at (0.77,1.05) {30};
    \node[font=\small] at (3.96,1.05) {23};
    \node[font=\small] at (4.66,1.05) {22};
    \node[font=\small] at (15.12,1.05) {0};


    \draw (0.48,-0.05) -- (0.48,-0.25);
    \draw (0.48,-0.25) -- (4.32,-0.25);
    \draw (4.32,-0.05) -- (4.32,-0.25);

    \node[font=\small] at (2.36,-0.65) {смещение порядка};

    \draw (4.32,-0.05) -- (4.32,-0.25);
    \draw (4.32,-0.25) -- (15.12,-0.25);
    \draw (15.12,-0.05) -- (15.12,-0.25);

    \node[font=\small] at (9.48,-0.65) {мантисса};

    \end{tikzpicture}
\end{figure}

\subsubsection{Представление числа В}

$B=0.0000 1100 1100 1100 1100 1100 1101_{2}$.

Получим представление числа в формате, зафиксировав запятую после первой единицы.

$B = (1.10011001100110011001101)_{2} \cdot 2^{-5}$

$X_{A} = -5 + 127 = 122_{10} = 01111010_{2}$

Дополним смещение порядка до 8 бит, запишем знаковый бит 0 и полученные значения
в таблицу, учитывая скрытый бит, и получим представление числа в формате Ф3.

\begin{figure}[H]
    \centering
    \begin{tikzpicture}[
        bit/.style={font=\large, minimum width=0.48cm,
                    minimum height=0.65cm, inner sep=0pt},
        every node/.style={inner sep=0pt}
    ]

    \def\bits{
        0,0,1,1,1,1,0,1,
        0,1,0,0,1,1,0,0,
        1,1,0,0,1,1,0,0,
        1,1,0,0,1,1,0,1
    }

    \draw (0,0) rectangle (15.36,0.7);
    \draw (0,0) rectangle (0.48,0.7);
    \draw (0,0) rectangle (4.32,0.7);

    \foreach \b [count=\i from 0] in \bits {
        \node[bit] at ({0.24+0.48*\i},0.35) {\b};
    }

    \node[font=\small] at (0.24,1.05) {31};
    \node[font=\small] at (0.77,1.05) {30};
    \node[font=\small] at (3.96,1.05) {23};
    \node[font=\small] at (4.66,1.05) {22};
    \node[font=\small] at (15.12,1.05) {0};


    \draw (0.48,-0.05) -- (0.48,-0.25);
    \draw (0.48,-0.25) -- (4.32,-0.25);
    \draw (4.32,-0.05) -- (4.32,-0.25);

    \node[font=\small] at (2.36,-0.65) {смещение порядка};

    \draw (4.32,-0.05) -- (4.32,-0.25);
    \draw (4.32,-0.25) -- (15.12,-0.25);
    \draw (15.12,-0.05) -- (15.12,-0.25);

    \node[font=\small] at (9.48,-0.65) {мантисса};

    \end{tikzpicture}
\end{figure}


\subsection{Получение чисел из формата Ф1}
\subsubsection{Получение числа Y из R}

Распишем число R в таблице~\ref{tab:r-in-f1}.

\begin{table}[H]
    \centering
    \caption{Двоичное представление числа R}
    \label{tab:r-in-f1}
    \begin{tabular}{|c|c|c|c|c|c|c|c|}
        \hline
		$C$ & $1$ & $E$ & $0$ & $8$ & $0$ & $0$ & $0$ \\
		\hline
		$1100$ & $0001$ & $1110$ & $0000$ & $1000$ & $0000$ & $0000$ & $0000$ \\
        \hline
    \end{tabular}
\end{table}

Расставим полученное число в формате Ф1.

\begin{figure}[H]
    \centering
    \begin{tikzpicture}[
        bit/.style={font=\large, minimum width=0.48cm,
                    minimum height=0.65cm, inner sep=0pt},
        every node/.style={inner sep=0pt}
    ]

    \def\bits{
        1,1,0,0,0,0,0,1,
        1,1,1,0,0,0,0,0,
        1,0,0,0,0,0,0,0,
        0,0,0,0,0,0,0,0
    }

    \draw (0,0) rectangle (15.36,0.7);
    \draw (0,0) rectangle (0.48,0.7);
    \draw (0,0) rectangle (3.84,0.7);

    \foreach \b [count=\i from 0] in \bits {
        \node[bit] at ({0.24+0.48*\i},0.35) {\b};
    }

    \node[font=\small] at (0.24,1.05) {0};
    \node[font=\small] at (0.72,1.05) {1};
    \node[font=\small] at (3.58,1.05) {7};
    \node[font=\small] at (4.06,1.05) {8};
    \node[font=\small] at (15.12,1.05) {31};


    \draw (0.48,-0.05) -- (0.48,-0.25);
    \draw (0.48,-0.25) -- (3.84,-0.25);
    \draw (3.84,-0.05) -- (3.84,-0.25);

    \node[font=\small] at (2.16,-0.65) {характеристика};

    \draw (3.84,-0.05) -- (3.84,-0.25);
    \draw (3.84,-0.25) -- (15.12,-0.25);
    \draw (15.12,-0.05) -- (15.12,-0.25);

    \node[font=\small] at (9.48,-0.65) {мантисса};

    \end{tikzpicture}
\end{figure}

Так как первый бит $= 1$, то число $Y < 0$. Получим порядок числа.

$P_{A}=X_{A}-d=(1000001)_{2}-64=65-64=1$

Значит, число имеет вид $Y = -(0.E08000)_{16} \cdot 16^1 = -E.08_{16}$.

Переведём число в десятичную систему:

$Y=-(14 \cdot 16^0 + 0 \cdot 16^{-1} + 8 \cdot 16^{-2})=-14.03125_{10}$

\subsubsection{Получение числа Z из S}

Распишем число S в таблице~\ref{tab:s-in-f1}.

\begin{table}[H]
    \centering
    \caption{Двоичное представление числа S}
    \label{tab:s-in-f1}
    \begin{tabular}{|c|c|c|c|c|c|c|c|}
        \hline
		$3$ & $E$ & $A$ & $0$ & $0$ & $0$ & $0$ & $0$ \\
		\hline
		$0011$ & $1110$ & $1010$ & $0000$ & $0000$ & $0000$ & $0000$ & $0000$ \\
        \hline
    \end{tabular}
\end{table}

Расставим полученное число в формате Ф1.

\begin{figure}[H]
    \centering
    \begin{tikzpicture}[
        bit/.style={font=\large, minimum width=0.48cm,
                    minimum height=0.65cm, inner sep=0pt},
        every node/.style={inner sep=0pt}
    ]

    \def\bits{
        0,0,1,1,1,1,1,0,
        1,0,1,0,0,0,0,0,
        0,0,0,0,0,0,0,0,
        0,0,0,0,0,0,0,0
    }

    \draw (0,0) rectangle (15.36,0.7);
    \draw (0,0) rectangle (0.48,0.7);
    \draw (0,0) rectangle (3.84,0.7);

    \foreach \b [count=\i from 0] in \bits {
        \node[bit] at ({0.24+0.48*\i},0.35) {\b};
    }

    \node[font=\small] at (0.24,1.05) {0};
    \node[font=\small] at (0.72,1.05) {1};
    \node[font=\small] at (3.58,1.05) {7};
    \node[font=\small] at (4.06,1.05) {8};
    \node[font=\small] at (15.12,1.05) {31};


    \draw (0.48,-0.05) -- (0.48,-0.25);
    \draw (0.48,-0.25) -- (3.84,-0.25);
    \draw (3.84,-0.05) -- (3.84,-0.25);

    \node[font=\small] at (2.16,-0.65) {характеристика};

    \draw (3.84,-0.05) -- (3.84,-0.25);
    \draw (3.84,-0.25) -- (15.12,-0.25);
    \draw (15.12,-0.05) -- (15.12,-0.25);

    \node[font=\small] at (9.48,-0.65) {мантисса};

    \end{tikzpicture}
\end{figure}

Так как первый бит $= 0$, то число $Z > 0$. Получим порядок числа.

$P_{A}=X_{A}-d=(0111110)_{2}-64=62-64=-2$

Значит, число имеет вид $Z = (0.A00000)_{16} \cdot 16^{-2} = 0.00A_{16}$.

Переведём число в десятичную систему:

$Z=(10 \cdot 16^{-3})=0.00244140625_{10}$


\subsection{Получение чисел из формата Ф2}
\subsubsection{Получение числа V из R}

Распишем число R в таблице~\ref{tab:r-in-f2}.

\begin{table}[H]
    \centering
    \caption{Двоичное представление числа R}
    \label{tab:r-in-f2}
    \begin{tabular}{|c|c|c|c|c|c|c|c|}
        \hline
		$C$ & $1$ & $E$ & $0$ & $8$ & $0$ & $0$ & $0$ \\
		\hline
		$1100$ & $0001$ & $1110$ & $0000$ & $1000$ & $0000$ & $0000$ & $0000$ \\
        \hline
    \end{tabular}
\end{table}

Расставим полученное число в формате Ф2.

\begin{figure}[H]
    \centering
    \begin{tikzpicture}[
        bit/.style={font=\large, minimum width=0.48cm,
                    minimum height=0.65cm, inner sep=0pt},
        every node/.style={inner sep=0pt}
    ]

    \def\bits{
        1,1,0,0,0,0,0,1,
        1,1,1,0,0,0,0,0,
        1,0,0,0,0,0,0,0,
        0,0,0,0,0,0,0,0
    }

    \draw (0,0) rectangle (15.36,0.7);
    \draw (0,0) rectangle (0.48,0.7);
    \draw (0,0) rectangle (4.32,0.7);

    \foreach \b [count=\i from 0] in \bits {
        \node[bit] at ({0.24+0.48*\i},0.35) {\b};
    }

    \node[font=\small] at (0.24,1.05) {31};
    \node[font=\small] at (0.77,1.05) {30};
    \node[font=\small] at (3.96,1.05) {23};
    \node[font=\small] at (4.66,1.05) {22};
    \node[font=\small] at (15.12,1.05) {0};


    \draw (0.48,-0.05) -- (0.48,-0.25);
    \draw (0.48,-0.25) -- (4.32,-0.25);
    \draw (4.32,-0.05) -- (4.32,-0.25);

    \node[font=\small] at (2.16,-0.65) {характеристика};

    \draw (4.32,-0.05) -- (4.32,-0.25);
    \draw (4.32,-0.25) -- (15.12,-0.25);
    \draw (15.12,-0.05) -- (15.12,-0.25);

    \node[font=\small] at (9.48,-0.65) {мантисса};

    \end{tikzpicture}
\end{figure}

Так как первый бит $= 1$, то число $V < 0$. Получим порядок числа.

$P_{A}=X_{A}-d=(10000011)_{2}-128=131-128=3$

Учитывая скрытый разряд, получим вид числа $V = -(0.111000001)_{2} \cdot 2^3 = -111.000001_{2}$.

Переведём число в десятичную систему:

$V=-(2^2+2^1+2^0+2^{-6})=-7.012625_{10}$

\subsubsection{Получение числа W из S}

Распишем число S в таблице~\ref{tab:s-in-f2}.

\begin{table}[H]
    \centering
    \caption{Двоичное представление числа S}
    \label{tab:s-in-f2}
    \begin{tabular}{|c|c|c|c|c|c|c|c|}
        \hline
		$3$ & $E$ & $A$ & $0$ & $0$ & $0$ & $0$ & $0$ \\
		\hline
		$0011$ & $1110$ & $1010$ & $0000$ & $0000$ & $0000$ & $0000$ & $0000$ \\
        \hline
    \end{tabular}
\end{table}

Расставим полученное число в формате Ф2.

\begin{figure}[H]
    \centering
    \begin{tikzpicture}[
        bit/.style={font=\large, minimum width=0.48cm,
                    minimum height=0.65cm, inner sep=0pt},
        every node/.style={inner sep=0pt}
    ]

    \def\bits{
        0,0,1,1,1,1,1,0,
        1,0,1,0,0,0,0,0,
        0,0,0,0,0,0,0,0,
        0,0,0,0,0,0,0,0
    }

    \draw (0,0) rectangle (15.36,0.7);
    \draw (0,0) rectangle (0.48,0.7);
    \draw (0,0) rectangle (4.32,0.7);

    \foreach \b [count=\i from 0] in \bits {
        \node[bit] at ({0.24+0.48*\i},0.35) {\b};
    }

    \node[font=\small] at (0.24,1.05) {31};
    \node[font=\small] at (0.77,1.05) {30};
    \node[font=\small] at (3.96,1.05) {23};
    \node[font=\small] at (4.66,1.05) {22};
    \node[font=\small] at (15.12,1.05) {0};


    \draw (0.48,-0.05) -- (0.48,-0.25);
    \draw (0.48,-0.25) -- (4.32,-0.25);
    \draw (4.32,-0.05) -- (4.32,-0.25);

    \node[font=\small] at (2.16,-0.65) {характеристика};

    \draw (4.32,-0.05) -- (4.32,-0.25);
    \draw (4.32,-0.25) -- (15.12,-0.25);
    \draw (15.12,-0.05) -- (15.12,-0.25);

    \node[font=\small] at (9.48,-0.65) {мантисса};

    \end{tikzpicture}
\end{figure}

Так как первый бит $= 0$, то число $W > 0$. Получим порядок числа.

$P_{A}=X_{A}-d=(01111101)_{2}-128=125-128=-3$

Учитывая скрытый разряд, получим вид числа $W = (0.101)_{2} \cdot 2^{-3} = 0.000101_{2}$.

Переведём число в десятичную систему:

$W=(2^{-4} + 2^{-6})=0.078125_{10}$


\subsection{Получение чисел из формата Ф3}
\subsubsection{Получение числа T из R}

Распишем число R в таблице~\ref{tab:r-in-f3}.

\begin{table}[H]
    \centering
    \caption{Двоичное представление числа R}
    \label{tab:r-in-f3}
    \begin{tabular}{|c|c|c|c|c|c|c|c|}
        \hline
		$C$ & $1$ & $E$ & $0$ & $8$ & $0$ & $0$ & $0$ \\
		\hline
		$1100$ & $0001$ & $1110$ & $0000$ & $1000$ & $0000$ & $0000$ & $0000$ \\
        \hline
    \end{tabular}
\end{table}

Расставим полученное число в формате Ф3.

\begin{figure}[H]
    \centering
    \begin{tikzpicture}[
        bit/.style={font=\large, minimum width=0.48cm,
                    minimum height=0.65cm, inner sep=0pt},
        every node/.style={inner sep=0pt}
    ]

    \def\bits{
        1,1,0,0,0,0,0,1,
        1,1,1,0,0,0,0,0,
        1,0,0,0,0,0,0,0,
        0,0,0,0,0,0,0,0
    }

    \draw (0,0) rectangle (15.36,0.7);
    \draw (0,0) rectangle (0.48,0.7);
    \draw (0,0) rectangle (4.32,0.7);

    \foreach \b [count=\i from 0] in \bits {
        \node[bit] at ({0.24+0.48*\i},0.35) {\b};
    }

    \node[font=\small] at (0.24,1.05) {31};
    \node[font=\small] at (0.77,1.05) {30};
    \node[font=\small] at (3.96,1.05) {23};
    \node[font=\small] at (4.66,1.05) {22};
    \node[font=\small] at (15.12,1.05) {0};


    \draw (0.48,-0.05) -- (0.48,-0.25);
    \draw (0.48,-0.25) -- (4.32,-0.25);
    \draw (4.32,-0.05) -- (4.32,-0.25);

    \node[font=\small] at (2.36,-0.65) {смещение порядка};

    \draw (4.32,-0.05) -- (4.32,-0.25);
    \draw (4.32,-0.25) -- (15.12,-0.25);
    \draw (15.12,-0.05) -- (15.12,-0.25);

    \node[font=\small] at (9.48,-0.65) {мантисса};

    \end{tikzpicture}
\end{figure}

Так как первый бит $= 1$, то число $T < 0$. Получим порядок числа.

$P_{A}=X_{A}-d=(10000011)_{2}-127=131-127=4$

Учитывая скрытый разряд, получим вид числа $T = -(1.11000001)_{2} \cdot 2^4 = -11100.0001_{2}$.

Переведём число в десятичную систему:

$T=-(2^4+2^3+2^2+2^{-4})=-28.0625_{10}$

\subsubsection{Получение числа Q из S}

Распишем число S в таблице~\ref{tab:s-in-f3}.

\begin{table}[H]
    \centering
    \caption{Двоичное представление числа S}
    \label{tab:s-in-f3}
    \begin{tabular}{|c|c|c|c|c|c|c|c|}
        \hline
		$3$ & $E$ & $A$ & $0$ & $0$ & $0$ & $0$ & $0$ \\
		\hline
		$0011$ & $1110$ & $1010$ & $0000$ & $0000$ & $0000$ & $0000$ & $0000$ \\
        \hline
    \end{tabular}
\end{table}

Расставим полученное число в формате Ф3.

\begin{figure}[H]
    \centering
    \begin{tikzpicture}[
        bit/.style={font=\large, minimum width=0.48cm,
                    minimum height=0.65cm, inner sep=0pt},
        every node/.style={inner sep=0pt}
    ]

    \def\bits{
        0,0,1,1,1,1,1,0,
        1,0,1,0,0,0,0,0,
        0,0,0,0,0,0,0,0,
        0,0,0,0,0,0,0,0
    }

    \draw (0,0) rectangle (15.36,0.7);
    \draw (0,0) rectangle (0.48,0.7);
    \draw (0,0) rectangle (4.32,0.7);

    \foreach \b [count=\i from 0] in \bits {
        \node[bit] at ({0.24+0.48*\i},0.35) {\b};
    }

    \node[font=\small] at (0.24,1.05) {31};
    \node[font=\small] at (0.77,1.05) {30};
    \node[font=\small] at (3.96,1.05) {23};
    \node[font=\small] at (4.66,1.05) {22};
    \node[font=\small] at (15.12,1.05) {0};


    \draw (0.48,-0.05) -- (0.48,-0.25);
    \draw (0.48,-0.25) -- (4.32,-0.25);
    \draw (4.32,-0.05) -- (4.32,-0.25);

    \node[font=\small] at (2.36,-0.65) {смещение порядка};

    \draw (4.32,-0.05) -- (4.32,-0.25);
    \draw (4.32,-0.25) -- (15.12,-0.25);
    \draw (15.12,-0.05) -- (15.12,-0.25);

    \node[font=\small] at (9.48,-0.65) {мантисса};

    \end{tikzpicture}
\end{figure}

Так как первый бит $= 0$, то число $Q > 0$. Получим порядок числа.

$P_{A}=X_{A}-d=(01111101)_{2}-127=125-127=-2$

Учитывая скрытый разряд, получим вид числа $Q = (1.01)_{2} \cdot 2^{-2} = 0.0101_{2}$.

Переведём число в десятичную систему:

$Q=(2^{-2} + 2^{-4})=0.3125_{10}$

\section{Сложение целых чисел}
\subsection{Сложение чисел A и В c различными знаками}

В таблице~\ref{tab:a-b-dop} представлены прямые и дополнительные коды
для чисел $A$ и $B$.

\begin{table}[H]
    \centering
    \caption{Прямые и дополнительные коды для A и В}
	\label{tab:a-b-dop}
    \begin{tabular}{|c|c|c|}
		\hline
		 & Прямой код & Дополнительный код \\
		\hline
		$A=109$ & $0110\ 1101$ & \makecell{$2^8-109$\\$1001\ 0011$}  \\
		\hline
		$B=17$ & $0001\ 0001$ & \makecell{$2^8-17$\\$1110\ 1111$} \\
		\hline
    \end{tabular}
\end{table}

\subsubsection{A+B}

\begin{table}[H]
    \centering
    \caption{Сложение А и В}
    \begin{tabular}{c|cc|cccccccp{1cm}ccp{0.5cm}cc}
		\small перенос & & & & & & & & $1$ & & & \multicolumn{2}{c}{знаковая} & \  & \multicolumn{2}{c}{беззнаковая}  \\
		\cline{1-10} \cline{12-13} \cline{15-16}
		$A_{\text{пр}}$ & & 0 & 1 & 1 & 0 & 1 & 1 & 0 & 1 & & \multirow{2}{*}{+} & 109 & \  & \multirow{2}{*}{+} & 109   \\
		$B_{\text{пр}}$ & & 0 & 0 & 0 & 1 & 0 & 0 & 0 & 1 & & & 17 & \  &  & 17 \\
		\cline{1-10} \cline{12-13} \cline{15-16}
		$C_{\text{пр}}$ & & 0 & 1 & 1 & 1 & 1 & 1 & 1 & 0 & & & 126 & \  &  & 126 \\
    \end{tabular}
\end{table}

\begin{table}[H]
    \centering
    \begin{tabular}{cccccc}
        CF=0; & ZF=0; & PF=1; & AF=0; & SF=0; & OF=0;
    \end{tabular}
\end{table}

\subsubsection{-A+B}

\begin{table}[H]
    \centering
    \caption{Сложение -А и В}
    \begin{tabular}{c|cc|cccccccp{1cm}ccp{0.5cm}cc}
		\small перенос & & & & $1$ & & & $1$ & $1$ & & & \multicolumn{2}{c}{знаковая} & \  & \multicolumn{2}{c}{беззнаковая}  \\
		\cline{1-10} \cline{12-13} \cline{15-16}
		$A_{\text{доп}}$ & & 1 & 0 & 0 & 1 & 0 & 0 & 1 & 1 & & \multirow{2}{*}{+} & -109 & \  & \multirow{2}{*}{+} & 147  \\
		$B_{\text{пр}}$ & & 0 & 0 & 0 & 1 & 0 & 0 & 0 & 1 & & & 17 & \  &  & 17 \\
		\cline{1-10} \cline{12-13} \cline{15-16}
		$C_{\text{доп}}$ & & 1 & 0 & 1 & 0 & 0 & 1 & 0 & 0 & & & & \  &  & 164 \\
		\cline{1-10} \cline{12-13} \cline{15-16}
		$C_{\text{пр}}$ & & 1 & 1 & 0 & 1 & 1 & 1 & 0 & 0 & & & -92 & \  &  & 
    \end{tabular}
\end{table}

\begin{table}[H]
    \centering
    \begin{tabular}{cccccc}
        CF=0; & ZF=0; & PF=0; & AF=0; & SF=1; & OF=0;
    \end{tabular}
\end{table}


\subsubsection{A+(-B)}

\begin{table}[H]
    \centering
    \caption{Сложение А и -В}
    \begin{tabular}{c|cc|cccccccp{1cm}ccp{0.5cm}cc}
		\small перенос & $1$ & $1$ & $1$ & & $1$ & $1$ & $1$ & $1$ & & & \multicolumn{2}{c}{знаковая} & \  & \multicolumn{2}{c}{беззнаковая}  \\
		\cline{1-10} \cline{12-13} \cline{15-16}
		$A_{\text{пр}}$ & & 0 & 1 & 1 & 0 & 1 & 1 & 0 & 1 & & \multirow{2}{*}{+} & 109 & \  & \multirow{2}{*}{+} & 109   \\
		$B_{\text{доп}}$ & & 1 & 1 & 1 & 0 & 1 & 1 & 1 & 1 & & & -17 & \  &  & 239 \\
		\cline{1-10} \cline{12-13} \cline{15-16}
		$C_{\text{пр}}$ & & 0 & 1 & 0 & 1 & 1 & 1 & 0 & 0 & & & 92 & \  &  & 92? \\
    \end{tabular}
\end{table}

\begin{table}[H]
    \centering
    \begin{tabular}{cccccc}
        CF=1; & ZF=0; & PF=1; & AF=1; & SF=0; & OF=0;
    \end{tabular}
\end{table}

Для беззнаковой интерпретации результат неверен из-за переноса старшего разряда,
величина которого 256. $109+239=348$, $348-256=92$. 


\subsubsection{-A+(-B)}

\begin{table}[H]
    \centering
    \caption{Сложение -А и -В}
    \begin{tabular}{c|cc|cccccccp{1cm}ccp{0.5cm}cc}
		\small перенос & $1$ & $1$ & $1$ & $1$ & $1$ & $1$ & $1$ & $1$ & & & \multicolumn{2}{c}{знаковая} & \  & \multicolumn{2}{c}{беззнаковая}  \\
		\cline{1-10} \cline{12-13} \cline{15-16}
		$A_{\text{доп}}$ & & 1 & 0 & 0 & 1 & 0 & 0 & 1 & 1 & & \multirow{2}{*}{+} & -109 & \  & \multirow{2}{*}{+} & 147\\
		$B_{\text{доп}}$ & & 1 & 1 & 1 & 0 & 1 & 1 & 1 & 1 & & & -17 & \  &  & 239 \\
		\cline{1-10} \cline{12-13} \cline{15-16}
		$C_{\text{доп}}$ & & 1 & 0 & 0 & 0 & 0 & 0 & 1 & 0 & & & & \  &  & 130? \\
		\cline{1-10} \cline{12-13} \cline{15-16}
		$C_{\text{пр}}$ & & 1 & 1 & 1 & 1 & 1 & 1 & 1 & 0 & & & -126 & \  &  & \\
    \end{tabular}
\end{table}

\begin{table}[H]
    \centering
    \begin{tabular}{cccccc}
        CF=1; & ZF=0; & PF=0; & AF=1; & SF=1; & OF=0;
    \end{tabular}
\end{table}

Для беззнаковой интерпретации результат неверен из-за переноса старшего разряда,
величина которого 256. $147+239=386$, $386-256=130$. 

\subsection{Подбор числа В}
Для того, чтобы получить переполнение, необходимо, чтобы $A+B>128$. Возьмём
$B=20_{10}=0001\ 0100_{2}$. Тогда $-B=2^8-20=236_{10}=1110\ 1100_{2}$.

\begin{table}[H]
    \centering
    \caption{Сложение А и В}
    \begin{tabular}{c|cc|cccccccp{1cm}ccp{0.5cm}cc}
		\small перенос & & $1$ & $1$ & $1$ & $1$ & $1$ & & & & & \multicolumn{2}{c}{знаковая} & \  & \multicolumn{2}{c}{беззнаковая}  \\
		\cline{1-10} \cline{12-13} \cline{15-16}
		$A_{\text{пр}}$ & & 0 & 1 & 1 & 0 & 1 & 1 & 0 & 1 & & \multirow{2}{*}{+} & 109 & \  & \multirow{2}{*}{+} & 109   \\
		$B_{\text{пр}}$ & & 0 & 0 & 0 & 1 & 0 & 1 & 0 & 0 & & & 20 & \  &  & 20 \\
		\cline{1-10} \cline{12-13} \cline{15-16}
		$C_{\text{доп}}$ & & 1 & 0 & 0 & 0 & 0 & 0 & 0 & 1 & & & & \  &  & 129 \\
		\cline{1-10} \cline{12-13} \cline{15-16}
		$C_{\text{пр}}$ & & 1 & 1 & 1 & 1 & 1 & 1 & 1 & 1 & & & -127? & \  &  & \\
    \end{tabular}
\end{table}

\begin{table}[H]
    \centering
    \begin{tabular}{cccccc}
        CF=0; & ZF=0; & PF=1; & AF=1; & SF=1; & OF=1;
    \end{tabular}
\end{table}

Для знакового сложения результат неверный из-за переполнения формата.

\begin{table}[H]
    \centering
    \caption{Сложение -А и -В}
    \begin{tabular}{c|cc|cccccccp{1cm}ccp{0.5cm}cc}
		\small перенос & $1$ &  &  &  &  &  & & & & & \multicolumn{2}{c}{знаковая} & \  & \multicolumn{2}{c}{беззнаковая}  \\
		\cline{1-10} \cline{12-13} \cline{15-16}
		$A_{\text{доп}}$ & & 1 & 0 & 0 & 1 & 0 & 0 & 1 & 1 & & \multirow{2}{*}{+} & -109 & \  & \multirow{2}{*}{+} & 147\\
		$B_{\text{доп}}$ & & 1 & 1 & 1 & 0 & 1 & 1 & 0 & 0 & & & -20 & \  &  & 236 \\
		\cline{1-10} \cline{12-13} \cline{15-16}
		$C_{\text{пр}}$ & & 0 & 1 & 1 & 1 & 1 & 1 & 1 & 1 & & & 127? & \  &  & 127?\\
    \end{tabular}
\end{table}

\begin{table}[H]
    \centering
    \begin{tabular}{cccccc}
        CF=1; & ZF=0; & PF=0; & AF=0; & SF=0; & OF=1;
    \end{tabular}
\end{table}

Для знакового сложения результат неверный из-за переполнения формата.
Для беззнакового сложения результат неверный из-за переноса старшего разряда. 

\subsection{Подбор числа А}
Для того чтобы при сложении возникало переполнение, а при сложении отрицательных
чисел результат был корректен, необходимо $A+B=128$. Значит, возьмём
$A=111_{10}=0110\ 1111_{2}$, тогда $-A=2^8-109=147_{10}=1001\ 0001_{2}$.



\begin{table}[H]
    \centering
    \caption{Сложение А и В}
    \begin{tabular}{c|cc|cccccccp{1cm}ccp{0.5cm}cc}
		\small перенос & & $1$ & $1$ & $1$ & $1$ & $1$ & $1$ & $1$ & & & \multicolumn{2}{c}{знаковая} & \  & \multicolumn{2}{c}{беззнаковая}  \\
		\cline{1-10} \cline{12-13} \cline{15-16}
		$A_{\text{пр}}$ & & 0 & 1 & 1 & 0 & 1 & 1 & 1 & 1 & & \multirow{2}{*}{+} & 111 & \  & \multirow{2}{*}{+} & 111   \\
		$B_{\text{пр}}$ & & 0 & 0 & 0 & 1 & 0 & 0 & 0 & 1 & & & 17 & \  &  & 17 \\
		\cline{1-10} \cline{12-13} \cline{15-16}
		$C_{\text{доп}}$ & & 1 & 0 & 0 & 0 & 0 & 0 & 0 & 0 & & & & \  &  & 128 \\
		\cline{1-10} \cline{12-13} \cline{15-16}
		$C_{\text{пр}}$ & & 1 & 0 & 0 & 0 & 0 & 0 & 0 & 0 & & & -128? & \  &  & \\
    \end{tabular}
\end{table}

\begin{table}[H]
    \centering
    \begin{tabular}{cccccc}
        CF=0; & ZF=0; & PF=0; & AF=1; & SF=1; & OF=1;
    \end{tabular}
\end{table}

Для знакового сложения результат неверен из-за переполнения формата.

\begin{table}[H]
    \centering
    \caption{Сложение -А и -В}
    \begin{tabular}{c|cc|cccccccp{1cm}ccp{0.5cm}cc}
		\small перенос & $1$ & $1$ & $1$ & $1$ & $1$ & $1$ & $1$ & $1$ & & & \multicolumn{2}{c}{знаковая} & \  & \multicolumn{2}{c}{беззнаковая}  \\
		\cline{1-10} \cline{12-13} \cline{15-16}
		$A_{\text{доп}}$ & & 1 & 0 & 0 & 1 & 0 & 0 & 0 & 1 & & \multirow{2}{*}{+} & -111 & \  & \multirow{2}{*}{+} & 145   \\
		$B_{\text{доп}}$ & & 1 & 1 & 1 & 0 & 1 & 1 & 1 & 1 & & & -17 & \  &  & 239\\
		\cline{1-10} \cline{12-13} \cline{15-16}
		$C_{\text{доп}}$ & & 1 & 0 & 0 & 0 & 0 & 0 & 0 & 0 & & & & \  &  & 128? \\
		\cline{1-10} \cline{12-13} \cline{15-16}
		$C_{\text{пр}}$ & & 1 & 0 & 0 & 0 & 0 & 0 & 0 & 0 & & & -128 & \  &  & \\
    \end{tabular}
\end{table}

\begin{table}[H]
    \centering
    \begin{tabular}{cccccc}
        CF=1; & ZF=0; & PF=0; & AF=1; & SF=1; & OF=0;
    \end{tabular}
\end{table}

Для беззнакового сложения результат неверен из-за переноса старшего разряда.

\section{Вычитание целых чисел}
\subsection{Вычитание чисел A и В c различными знаками}

В таблице~\ref{tab:a-b-dop2} представлены прямые и дополнительные коды
для чисел $A$ и $B$.

\begin{table}[H]
    \centering
    \caption{Прямые и дополнительные коды для A и В}
	\label{tab:a-b-dop2}
    \begin{tabular}{|c|c|c|}
		\hline
		 & Прямой код & Дополнительный код \\
		\hline
		$A=109$ & $0110\ 1101$ & \makecell{$2^8-109$\\$1001\ 0011$}  \\
		\hline
		$B=17$ & $0001\ 0001$ & \makecell{$2^8-17$\\$1110\ 1111$} \\
		\hline
    \end{tabular}
\end{table}

\subsubsection{A-B}

\begin{table}[H]
    \centering
    \caption{Вычитание А и В}
    \begin{tabular}{c|cc|cccccccp{1cm}ccp{0.5cm}cc}
		\small заём & & & & $\bullet$ & $2$ & & & & & & \multicolumn{2}{c}{знаковая} & \  & \multicolumn{2}{c}{беззнаковая}  \\
		\cline{1-10} \cline{12-13} \cline{15-16}
		$A_{\text{пр}}$ & & 0 & 1 & 1 & 0 & 1 & 1 & 0 & 1 & & \multirow{2}{*}{-} & 109 & \  & \multirow{2}{*}{-} & 109   \\
		$B_{\text{пр}}$ & & 0 & 0 & 0 & 1 & 0 & 0 & 0 & 1 & & & 17 & \  &  & 17 \\
		\cline{1-10} \cline{12-13} \cline{15-16}
		$C_{\text{пр}}$ & & 0 & 1 & 0 & 1 & 1 & 1 & 0 & 0 & & & 92 & \  &  & 92 \\
    \end{tabular}
\end{table}

\begin{table}[H]
    \centering
    \begin{tabular}{cccccc}
        CF=0; & ZF=0; & PF=1; & AF=0; & SF=0; & OF=0;
    \end{tabular}
\end{table}

\subsubsection{-A-B}

\begin{table}[H]
    \centering
    \caption{Вычитание -А и В}
    \begin{tabular}{c|cc|cccccccp{1cm}ccp{0.5cm}cc}
		\small заём & & & & & & & & & & & \multicolumn{2}{c}{знаковая} & \  & \multicolumn{2}{c}{беззнаковая}  \\
		\cline{1-10} \cline{12-13} \cline{15-16}
		$A_{\text{доп}}$ & & 1 & 0 & 0 & 1 & 0 & 0 & 1 & 1 & & \multirow{2}{*}{-} & -109 & \  & \multirow{2}{*}{-} & 147  \\
		$B_{\text{пр}}$ & & 0 & 0 & 0 & 1 & 0 & 0 & 0 & 1 & & & 17 & \  &  & 17 \\
		\cline{1-10} \cline{12-13} \cline{15-16}
		$C_{\text{доп}}$ & & 1 & 0 & 0 & 0 & 0 & 0 & 1 & 0 & & & & \  &  & 130 \\
		\cline{1-10} \cline{12-13} \cline{15-16}
		$C_{\text{пр}}$ & & 1 & 1 & 1 & 1 & 1 & 1 & 1 & 0 & & & -126 & \  &  & 
    \end{tabular}
\end{table}

\begin{table}[H]
    \centering
    \begin{tabular}{cccccc}
        CF=0; & ZF=0; & PF=0; & AF=0; & SF=1; & OF=0;
    \end{tabular}
\end{table}


\subsubsection{A-(-B)}

\begin{table}[H]
    \centering
    \caption{Вычитание А и -В}
    \begin{tabular}{c|cc|cccccccp{1cm}ccp{0.5cm}cc}
		\small заём & $\bullet$ & $1$ & $1$ & $1$ & $1$ & $1$ & $1$ & $2$ & & & \multicolumn{2}{c}{знаковая} & \  & \multicolumn{2}{c}{беззнаковая}  \\
		\cline{1-10} \cline{12-13} \cline{15-16}
		$A_{\text{пр}}$ & & 0 & 1 & 1 & 0 & 1 & 1 & 0 & 1 & & \multirow{2}{*}{-} & 109 & \  & \multirow{2}{*}{-} & 109   \\
		$B_{\text{доп}}$ & & 1 & 1 & 1 & 0 & 1 & 1 & 1 & 1 & & & -17 & \  &  & 239 \\
		\cline{1-10} \cline{12-13} \cline{15-16}
		$C_{\text{пр}}$ & & 0 & 1 & 1 & 1 & 1 & 1 & 1 & 0 & & & 126 & \  &  & 126? \\
    \end{tabular}
\end{table}

\begin{table}[H]
    \centering
    \begin{tabular}{cccccc}
        CF=1; & ZF=0; & PF=1; & AF=1; & SF=0; & OF=0;
    \end{tabular}
\end{table}

Для беззнаковой интерпретации результат неверен из-за заёма разряда за пределами формата.

\subsubsection{-A-(-B)}

\begin{table}[H]
    \centering
    \caption{Вычитание -А и -В}
    \begin{tabular}{c|cc|cccccccp{1cm}ccp{0.5cm}cc}
		\small заём & $\bullet$ & $1$ & $1$ & $2$ & $\bullet$ & $1$ & $2$ & & & & \multicolumn{2}{c}{знаковая} & \  & \multicolumn{2}{c}{беззнаковая}  \\
		\cline{1-10} \cline{12-13} \cline{15-16}
		$A_{\text{доп}}$ & & 1 & 0 & 0 & 1 & 0 & 0 & 1 & 1 & & \multirow{2}{*}{-} & -109 & \  & \multirow{2}{*}{-} & 147\\
		$B_{\text{доп}}$ & & 1 & 1 & 1 & 0 & 1 & 1 & 1 & 1 & & & -17 & \  &  & 239 \\
		\cline{1-10} \cline{12-13} \cline{15-16}
		$C_{\text{доп}}$ & & 1 & 0 & 1 & 0 & 0 & 1 & 0 & 0 & & & & \  &  & 164? \\
		\cline{1-10} \cline{12-13} \cline{15-16}
		$C_{\text{пр}}$ & & 1 & 1 & 0 & 1 & 1 & 1 & 0 & 0 & & & -92 & \  &  & \\
    \end{tabular}
\end{table}

\begin{table}[H]
    \centering
    \begin{tabular}{cccccc}
        CF=1; & ZF=0; & PF=0; & AF=1; & SF=1; & OF=0;
    \end{tabular}
\end{table}

Для беззнаковой интерпретации результат неверен из-за заёма разряда за пределами формата.

\subsection{Подбор числа В}
Для того, чтобы получить переполнение, необходимо, чтобы $A+B>128$. Возьмём
$B=20_{10}=0001\ 0100_{2}$. Тогда $-B=2^8-20=236_{10}=1110\ 1100_{2}$.

\begin{table}[H]
    \centering
    \caption{Вычитание -А и В}
    \begin{tabular}{c|cc|cccccccp{1cm}ccp{0.5cm}cc}
		\small заём & & $\bullet$ & $1$ & $1$ & $1$ & $1$ & $2$ & & & & \multicolumn{2}{c}{знаковая} & \  & \multicolumn{2}{c}{беззнаковая}  \\
		\cline{1-10} \cline{12-13} \cline{15-16}
		$A_{\text{доп}}$ & & 1 & 0 & 0 & 1 & 0 & 0 & 1 & 1 & & \multirow{2}{*}{-} & -109 & \  & \multirow{2}{*}{-} & 147   \\
		$B_{\text{пр}}$ & & 0 & 0 & 0 & 1 & 0 & 1 & 0 & 0 & & & 20 & \  &  & 20 \\
		\cline{1-10} \cline{12-13} \cline{15-16}
		$C_{\text{пр}}$ & & 0 & 1 & 1 & 1 & 1 & 1 & 1 & 1 & & & 127? & \  &  & 127 \\
    \end{tabular}
\end{table}

\begin{table}[H]
    \centering
    \begin{tabular}{cccccc}
        CF=0; & ZF=0; & PF=1; & AF=1; & SF=1; & OF=1;
    \end{tabular}
\end{table}

Для знакового вычитания результат неверный из-за переполнения формата.

\begin{table}[H]
    \centering
    \caption{Вычитание А и -В}
    \begin{tabular}{c|cc|cccccccp{1cm}ccp{0.5cm}cc}
		\small заём & $\bullet$ & $2$ &  &  &  &  & & & & & \multicolumn{2}{c}{знаковая} & \  & \multicolumn{2}{c}{беззнаковая}  \\
		\cline{1-10} \cline{12-13} \cline{15-16}
		$A_{\text{пр}}$ & & 0 & 1 & 1 & 0 & 1 & 1 & 0 & 1 & & \multirow{2}{*}{-} & 109 & \  & \multirow{2}{*}{-} & 109\\
		$B_{\text{доп}}$ & & 1 & 1 & 1 & 0 & 1 & 1 & 0 & 0 & & & -20 & \  &  & 236 \\
		\cline{1-10} \cline{12-13} \cline{15-16}
		$C_{\text{доп}}$ & & 1 & 0 & 0 & 0 & 0 & 0 & 0 & 1 & & & & \  &  & 129?\\
		\cline{1-10} \cline{12-13} \cline{15-16}
		$C_{\text{пр}}$ & & 1 & 1 & 1 & 1 & 1 & 1 & 1 & 1 & & & -127? & \  &  & \\
    \end{tabular}
\end{table}

\begin{table}[H]
    \centering
    \begin{tabular}{cccccc}
        CF=1; & ZF=0; & PF=0; & AF=0; & SF=0; & OF=1;
    \end{tabular}
\end{table}

Для знакового вычитания результат неверный из-за переполнения формата.
Для беззнакового вычитания результат неверный из-за заёма разряда за пределами формата. 

\subsection{Подбор числа А}
Для того чтобы при $A-(-B)$ давало переполнение, а при $-A-B$
результат был корректен, необходимо $A+B=128$. Значит, возьмём
$A=111_{10}=0110\ 1111_{2}$, тогда $-A=2^8-109=147_{10}=1001\ 0001_{2}$.

\begin{table}[H]
    \centering
    \caption{Вычитание А и -В}
    \begin{tabular}{c|cc|cccccccp{1cm}ccp{0.5cm}cc}
		\small заём & $\bullet$ & $2$ & & & & & & & & & \multicolumn{2}{c}{знаковая} & \  & \multicolumn{2}{c}{беззнаковая}  \\
		\cline{1-10} \cline{12-13} \cline{15-16}
		$A_{\text{пр}}$ & & 0 & 1 & 1 & 0 & 1 & 1 & 1 & 1 & & \multirow{2}{*}{-} & 111 & \  & \multirow{2}{*}{-} & 111   \\
		$B_{\text{доп}}$ & & 1 & 1 & 1 & 0 & 1 & 1 & 1 & 1 & & & -17 & \  &  & 239\\
		\cline{1-10} \cline{12-13} \cline{15-16}
		$C_{\text{доп}}$ & & 1 & 0 & 0 & 0 & 0 & 0 & 0 & 0 & & & & \  &  & -128? \\
		\cline{1-10} \cline{12-13} \cline{15-16}
		$C_{\text{пр}}$ & & 1 & 0 & 0 & 0 & 0 & 0 & 0 & 0 & & & -128? & \  &  & \\
    \end{tabular}
\end{table}

\begin{table}[H]
    \centering
    \begin{tabular}{cccccc}
        CF=1; & ZF=0; & PF=0; & AF=0; & SF=1; & OF=1;
    \end{tabular}
\end{table}

Для знакового вычитания результат неверен из-за переполнения формата.
Для беззнакового вычитания результат неверен из-за заёма разряда за пределами формата.

\begin{table}[H]
    \centering
    \caption{Вычитание -А и В}
    \begin{tabular}{c|cc|cccccccp{1cm}ccp{0.5cm}cc}
		\small заём & & & & & & & & & & & \multicolumn{2}{c}{знаковая} & \  & \multicolumn{2}{c}{беззнаковая}  \\
		\cline{1-10} \cline{12-13} \cline{15-16}
		$A_{\text{доп}}$ & & 1 & 0 & 0 & 1 & 0 & 0 & 0 & 1 & & \multirow{2}{*}{-} & -111 & \  & \multirow{2}{*}{-} & 145   \\
		$B_{\text{пр}}$ & & 0 & 0 & 0 & 1 & 0 & 0 & 0 & 1 & & & 17 & \  &  & 17\\
		\cline{1-10} \cline{12-13} \cline{15-16}
		$C_{\text{доп}}$ & & 1 & 0 & 0 & 0 & 0 & 0 & 0 & 0 & & & & \  &  & 128 \\
		\cline{1-10} \cline{12-13} \cline{15-16}
		$C_{\text{пр}}$ & & 1 & 0 & 0 & 0 & 0 & 0 & 0 & 0 & & & -128 & \  &  & \\
    \end{tabular}
\end{table}

\begin{table}[H]
    \centering
    \begin{tabular}{cccccc}
        CF=0; & ZF=0; & PF=0; & AF=0; & SF=1; & OF=0;
    \end{tabular}
\end{table}


\section*{Заключение}
\addcontentsline{toc}{section}{Заключение}

В ходе выполнения работы были рассмотрены различные формы представления чисел
с фиксированной запятой и с плавающей запятой, рассмотрены примеры сложения
и вычитания целых чисел, а также их поведение при переносе, заёме и переполнении.



\end{document}
